\label{chap:introduction}

\section{Clinical Background and Motivation}
Type 2 Diabetes is at the forefront of metabolic issues presenting a major 21st-century public health challenge. According to the International Diabetes Federation (IDF) Diabetes Atlas, over 537 million adults worldwide are currently living with diabetes, a figure projected to increase to 783 million by 2045 \cite{idf2021}. The chronic pathophysiology of Type 2 Diabetes involves a progressive decline in pancreatic beta-cell function, systemic insulin resistance in peripheral tissues (skeletal muscle and adipose tissue), and impaired hepatic glucose suppression.

While for many years clinical care has focused on fasting plasma glucose (FPG) and glycated hemoglobin (HbA1c) as the main tools in diagnosis, we now see from large scale epidemiological and physiological studies that Postprandial Glycemic Response (PPGR)—which is the degree, length and speed of the blood sugar rise after a meal—is a large player in glycemic fluctuation and cardiovascular issues \cite{brownlee2005pathobiology}. We note also that repeated postprandial glucose spikes above what is considered the physiologic threshold ($180\text{ mg/dL}$) cause acute endothelial damage, increase the production of reactive oxygen species by mitochondria, set off the production of advanced glycation end products (AGEs), and in turn bring about microvascular complications which include diabetic nephropathy, retinopathy and peripheral neuropathy \cite{brownlee2005pathobiology}.

\section{The Dilemma of Inter-Individual Glycemic Variability}
For what has been a long time in diabetology standard nutritional counseling has put forth a one size fits all approach. Diabetics and pre-diabetics are told to follow the same carbohydrate counts, universal glycemic index tables, and broad caloric restriction. But recent clinical trials by Zeevi et al. \cite{zeevi2015personalized}, Hall et al. \cite{hall2018glucotypes}, and Berry et al. \cite{berry2020human} challenged the base of these broad-stroke approaches.

These studies report that in many cases we see great differences between what one person's blood sugar does versus another's after they eat the same meal with the exact same amount of carbs. In the research done by Zeevi et al. \cite{zeevi2015personalized}, some study participants had very large postprandial spikes ($>200\text{ mg/dL}$) after eating white bread, while others had very little response to that bread but had large spikes after eating fruit high in fructose. This large inter-individual variation is a result of complex non-linear interactions across several physiological systems:

\begin{enumerate}
    \item \textbf{Baseline Insulin Sensitivity and Secretion:} Quantified clinically using Homeostatic Model Assessment of Insulin Resistance (HOMA-IR), fasting insulin levels, and hemoglobin HbA1c.
    \item \textbf{Real-Time Activity and Motion Dynamics:} Muscle movement activates insulin-independent GLUT4 glucose transporters and directly changes how fast glucose clears from the blood.
    \item \textbf{Pre-Meal Glycemic Velocity:} The physiological state of glucose momentum (rising, falling, or steady) at the exact moment a meal is eaten.
    \item \textbf{Gut Microbiome Taxonomic Composition:} Thousands of bacterial species in the gut ferment dietary fiber into short-chain fatty acids (SCFAs), regulate GLP-1 hormone release, and influence systemic inflammation.
\end{enumerate}

\section{Limitations of Existing Continuous Glucose Monitoring Approaches}
Wearable Continuous Glucose Monitors (CGMs), such as the Abbott FreeStyle Libre and Dexcom G6, have made tracking daily blood sugar much easier. These sensors measure interstitial glucose every 1 to 5 minutes so users can follow their trends throughout the day. Even though they provide continuous data, current monitoring setups have three clear practical limitations:

\begin{enumerate}
    \item \textbf{Reactive Rather than Proactive Alerts:} Commercial CGM companion apps only notify users after a spike or drop has already taken place (such as above $180\text{ mg/dL}$ or below $70\text{ mg/dL}$). Once an alert sounds, the food has already been eaten and digested. The user cannot undo the meal, which makes the alert purely retrospective rather than preventive.
    \item \textbf{Isolated Glucose Tracking Without Context:} Most monitoring tools look at glucose values in isolation. They do not easily connect these sensor readings with actual meal nutrients (carbs, fiber, protein, and fat), smartwatch activity levels (heart rate and METs), standard blood lab tests, or gut microbiome profiles.
    \item \textbf{Handling High-Dimensional Gut Data:} Gut microbiome reports contain presence-absence data for thousands of bacteria. Because study cohorts usually have fewer than a hundred participants, existing pipelines either skip this data completely or use simple linear PCA. Linear PCA assumes continuous bell-curve distributions, which does not fit sparse zero-one bacterial data and misses how bacteria interact in groups.
\end{enumerate}

\section{Problem Statement and Scope}
Because current CGMs only warn users after blood sugar has already gone up, there is a clear need for a personalized digital twin tool. Such a system should predict full post-meal glucose curves ahead of time and allow patients to simulate different food choices before eating.

The scope of this thesis encompasses:
\begin{itemize}
    \item Cleaning and aligning sensor logs from 45 participants (1,633 logged meal events) in the CGMacros study \cite{mendes2024cgmacros}, covering healthy individuals, pre-diabetic patients, and individuals with Type 2 Diabetes.
    \item Designing an unsupervised representation learning framework tailored for sparse, high-dimensional binary gut microbiome taxonomy.
    \item Formulating an attention-based deep learning architecture to predict continuous 120-minute postprandial glucose trajectories on completely unseen human subjects.
    \item Formulating a structural causal simulation engine enabling patients to evaluate dietary substitutions and exercise timing before eating.
    \item Constructing and deploying a responsive, mobile-first Progressive Web App (PWA) running real-time inferences for end-users.
\end{itemize}

\section{Research Objectives}
To address the challenges described above, this project focuses on five main objectives:
\begin{enumerate}
    \item \textbf{Multi-Modal Data Pipeline and Sensor Alignment:} Build an automated preprocessing workflow that aligns dual CGM recordings (Libre and Dexcom), Fitbit activity logs (heart rate and METs), meal macronutrient grams, standard clinical lab values, and 1,979 gut bacterial species into unified time-series samples.
    \item \textbf{Microbiome Feature Extraction Using a $\beta$-VAE:} Train an unsupervised $\beta$-Variational Autoencoder with input noise masking and annealing to compress 1,979 binary gut species into 8 continuous latent variables ($\mathbf{z} \in \mathbb{R}^8$), capturing useful microbial patterns without overfitting the small cohort.
    \item \textbf{Continuous Glucose Curve Forecasting:} Construct a hybrid deep learning model called TransformerForecaster. It uses self-attention over 60 minutes of pre-meal sensor trends, gated residual blocks for patient and meal traits, a residual delta target ($\Delta G_t = G_t - G_0$) to avoid baseline shifts, and a combined loss function that penalizes peak and slope errors.
    \item \textbf{Model Evaluation and Clinical Safety Testing:} Test model performance using 5-fold cross-validation grouped by patient so test subjects are strictly unseen. Compare the results against baseline MLP, LSTM, and GRU models, and verify clinical safety using Clarke Error Grid standards.
    \item \textbf{What-If Simulation and Mobile Interface:} Create a causal simulation module that ranks different lifestyle options using a Clinical Utility Score, and build a simple mobile web app so users can test meal adjustments and walks directly on their phones.
\end{enumerate}

\section{Key Contributions of the Thesis}
The primary contributions of this thesis are summarized as follows:
\begin{itemize}
    \item \textbf{An End-to-End Digital Twin System:} We built GlycoTwin, a pipeline that combines wearable streams, dietary macronutrients, blood test markers, and gut microbiome features to predict postprandial glucose curves.
    \item \textbf{Non-Linear Gut Embeddings:} We showed that an 8-variable $\beta$-VAE latent representation captures gut taxonomic data better than standard linear PCA, reaching strong reconstruction scores and pointing out relevant bacteria such as \emph{Oxalobacter formigenes} and \emph{Furfurilactobacillus rossiae}.
    \item \textbf{Accurate Continuous Curve Predictions:} When evaluated on completely new subjects, TransformerForecaster achieved an average curve error (MAE) of \textbf{19.37 mg/dL}, RMSE of \textbf{28.14 mg/dL}, and peak error of \textbf{24.73 mg/dL}, outperforming recurrent network baselines.
    \item \textbf{Clinical Accuracy and Safety Verification:} Based on Clarke Error Grid analysis, \textbf{98.47\% of predictions fell within safe zones (Zones A and B)}, with no predictions entering high-risk error regions (Zones D or E).
    \item \textbf{Actionable Proactive Decision Support:} We demonstrate how in silico counterfactual simulation can transform reactive diabetes tracking into proactive lifestyle therapy, blunting acute postprandial rises by over 85\% in high-risk individuals.
\end{itemize}

\section{Organization of the Thesis}
The remainder of this thesis is structured as follows:
\begin{itemize}
    \item \textbf{Chapter \ref{chap:literature_review}: Literature Review} reviews past work on mathematical glucose models, machine learning approaches, deep recurrent and transformer networks, gut microbiome studies, and digital twin methods.
    \item \textbf{Chapter \ref{chap:dataset_preprocessing}: Multi-Modal Dataset and Preprocessing Pipeline} describes the CGMacros cohort, explains the five types of collected data, and details sensor alignment and meal extraction steps.
    \item \textbf{Chapter \ref{chap:microbiome_vae}: Microbiome $\beta$-Variational Autoencoder} derives the mathematical formulation of our deep $\beta$-VAE, stochastic taxa-masking denoising, $\beta$-annealing schedule, and latent manifold discovery.
    \item \textbf{Chapter \ref{chap:transformer_forecaster}: TransformerForecaster Deep Learning Model} details the Multi-Head Self-Attention temporal encoder, Gated Residual Networks, Residual Delta modeling, and the Composite Excursion Loss function.
    \item \textbf{Chapter \ref{chap:results_benchmarks}: Experimental Results and Discussion} discusses the 5-fold cross-validation performance, compares different network architectures, and presents the Clarke Error Grid safety results.
    \item \textbf{Chapter \ref{chap:causal_simulator_app}: Causal Counterfactual Simulation and Mobile Application} explains the causal what-if simulation tool, the Clinical Utility Score formulation, and how the mobile web interface was implemented.
    \item \textbf{Chapter \ref{chap:conclusion}: Conclusion and Scope for Further Research} summarizes the project results, notes practical limitations, and suggests areas for future improvement.
\end{itemize}
